\section{Introduction}

Language-model agents are often adapted by placing a textual \emph{skill}---procedural instructions, tool conventions, and lessons distilled from prior successes---before each task.  This is attractive because the base model remains frozen, but long instructions consume context, attention compute, and cache memory on every invocation.  Prompt tuning offers the complementary idea of optimizing a small sequence of continuous embeddings while freezing the model~\citep{lester2021prompt,li2021prefix}.  The recent \softskill framework connects these ideas: it collects successful agent trajectories under a textual skill and trains a short soft prefix by next-token prediction on the entire trajectory~\citep{tao2026softskill}.

Full-trajectory supervision implicitly assumes that every token is an equally useful target for the prefix.  In a long code-producing trajectory, however, much of the syntax, boilerplate, and local lexical choice may already be likely under the frozen model without the skill.  A short prefix then spends capacity reproducing tokens that do not encode the skill's behavioral effect.  Worse, moving all positions toward one successful trajectory may perturb otherwise competent no-skill behavior.  These observations motivate a different target: compress the \emph{conditional change} caused by the skill, rather than the entire observed sequence.

We introduce \method, illustrated in Figure~\ref{fig:overview}.  For a task and a successful trajectory, we evaluate the same frozen model on identical gold prefixes with and without the hard skill.  This paired construction gives aligned next-token distributions without comparing two free-running generations.  We rank positions using both (i) how much the skill raises the probability of the realized gold token and (ii) how much it changes the full output distribution.  A fixed-budget selector retains only the strongest positions.  The student is the unchanged frozen model plus a length-8 soft prefix; it receives cross-entropy on selected gold tokens and a KL preservation loss toward the no-skill model at a matched set of unselected positions.

Our main evidence comes from SpreadsheetBench Verified~\citep{ma2024spreadsheetbench}, where success requires generating executable Python that correctly modifies a workbook.  The benchmark is deliberately challenging for teacher-forced compression: a single early generation error can change the program structure and invalidate all later gold contexts.  We use 61 successful trajectories generated for 80 training tasks, select on training trajectories only, tune on 40 validation tasks, and touch the 280-task test set once after freezing the configuration.  We make three contributions:

\begin{itemize}
    \item We formulate \emph{behavioral-delta localization} for prompt compression, using a paired hard-skill/no-skill teacher on the same successful prefixes.  Positive gain is highly concentrated: 5\% of positions account for 84.05\% of the total positive gain.
    \item We show that a selective length-8 prefix with an explicit no-skill preservation term improves matched test success over full-trajectory \softskill by 5.71 points (101/280 vs.\ 85/280).  Paired execution outcomes make this improvement statistically significant, although the method does not significantly beat the no-skill base.
    \item We report systematic negative results and failure diagnostics.  Wider local windows degrade validation success; sequential updates of overlapping prefix blocks do not behave like independent boosting; and teacher-forced residual closure correlates imperfectly with free-generation success.  These findings isolate exposure and state-distribution mismatch as the next bottleneck.
\end{itemize}

The intended claim is therefore precise: selection improves behavioral compression relative to the full-trajectory soft-prefix baseline in this setting.  We do not claim that an eight-token prefix has yet internalized the full textual skill, nor that teacher-forced token metrics alone predict agent success.

\begin{figure*}[t]
\centering
\begin{tikzpicture}[
  font=\footnotesize,
  node distance=6mm and 3mm,
  box/.style={draw,rounded corners=2pt,align=center,minimum height=10mm,inner sep=2.5pt},
  arr/.style={-{Latex[length=2mm]},thick},
  teacher/.style={box,fill=skillblue!8,draw=skillblue},
  student/.style={box,fill=skillgreen!8,draw=skillgreen},
  score/.style={box,fill=skillorange!8,draw=skillorange}
]
\node[box,text width=21mm] (data) {task $x$ + successful trajectory $y_{1:T}$};
\node[teacher,right=of data,text width=43mm] (paired) {same frozen LM, identical gold context\\hard skill: $q_t=p_\theta(\cdot\mid \skill,x,y_{<t})$\\no skill: $b_t=p_\theta(\cdot\mid x,y_{<t})$};
\node[score,right=of paired,text width=33mm] (locator) {behavioral-delta locator\\$g_t=[\log q_t(y_t)-\log b_t(y_t)]_+$\\$c_t=g_t\,\mathrm{JS}(q_t,b_t)$};
\node[score,right=of locator,text width=21mm] (select) {top-$\rho$ core $C$\\matched preserve set $U$};
\node[student,below=of locator,xshift=-13mm,text width=51mm] (train) {freeze $\theta$; train only $\prefix\in\mathbb{R}^{L\times d}$\\$\mathcal{L}_{\rm CE}(C)+\lambda\mathcal{L}_{\rm preserve}(U)$};
\node[student,right=of train,text width=30mm] (eval) {soft prefix + task\\free generation\\execution score};
\draw[arr] (data) -- (paired);
\draw[arr] (paired) -- (locator);
\draw[arr] (locator) -- (select);
\draw[arr] (select.south) |- (train.north);
\draw[arr] (train) -- (eval);
\end{tikzpicture}
\caption{\method separates \emph{where the hard skill changes behavior} from \emph{what the soft prefix must learn}.  The full-vocabulary distributions are used to locate a sparse core; current main training uses hard next-token targets on that core and a no-skill Top-64-plus-residual KL on unselected positions.  Only the soft-prefix embeddings are updated.}
\label{fig:overview}
\end{figure*}

